\documentclass[11pt,a4paper]{article}

% ============================================================================
%  Teoría del FD3D dinámico (Madariaga) tal como está implementado en el
%  código legacy de KDEllipsPy  (legacy/Dynamic_inversion/)
%  Autor de la nota: documentación interna KDEllipsPy
% ============================================================================

\usepackage[utf8]{inputenc}
\usepackage[T1]{fontenc}
\usepackage[spanish,es-noindentfirst]{babel}
\usepackage{amsmath,amssymb}
\usepackage{geometry}
\geometry{margin=2.4cm}
\usepackage{booktabs}
\usepackage{array}
\usepackage{xcolor}
\usepackage{listings}
\usepackage{hyperref}
\hypersetup{colorlinks=true,linkcolor=blue!55!black,urlcolor=blue!55!black}

\lstset{
  basicstyle=\ttfamily\small,
  breaklines=true,
  frame=single,
  backgroundcolor=\color{gray!8},
  columns=fullflexible,
}

\newcommand{\dd}{\mathrm{d}}
\newcommand{\pd}[2]{\frac{\partial #1}{\partial #2}}
\newcommand{\half}{\tfrac{1}{2}}

\title{\textbf{Teoría del solver dinámico FD3D (Madariaga)}\\
       \large Tal como está implementado en \texttt{legacy/Dynamic\_inversion/} de KDEllipsPy}
\author{Documentación interna --- KDEllipsPy (parte \emph{Dynamic})}
\date{\today}

\begin{document}
\maketitle

\begin{abstract}
Este documento describe, ecuación por ecuación y subrutina por subrutina, la
física que resuelve el código dinámico legacy de KDEllipsPy: un solver de
\textbf{ruptura espontánea 3D} por diferencias finitas en una malla escalonada
velocidad--esfuerzo (Madariaga 1976; Virieux 1986), con fricción
\emph{slip-weakening} aplicada sobre la componente $\sigma_{yz}$ de una falla
plana, acoplado a las funciones de Green 1D de AXITRA (número de onda discreto)
para sintetizar sismogramas. El objetivo es servir de base para reimplementar el
\emph{forward} en Python. Todas las afirmaciones están referidas a los fuentes
en \texttt{Source/fd3d\_subs/}.
\end{abstract}

\tableofcontents

% ===========================================================================
\section{Panorama: la cadena del problema directo}
% ===========================================================================

Una sola evaluación \emph{forward} (subrutina \texttt{forward\_modelling.f})
toma un vector de 10 parámetros del modelo \texttt{rmodel} y produce sismogramas
sintéticos, en esta cadena:

\begin{enumerate}
  \item \texttt{mkstress} --- construye el campo de \textbf{esfuerzo inicial}
        $\sigma^0(i,k)$ sobre la falla: una aspereza elíptica de alto esfuerzo
        más un círculo de nucleación.
  \item \texttt{mkpeak} --- construye el campo de \textbf{esfuerzo de cedencia}
        (\emph{peak/yield stress}) $\tau_p(i,k)$.
  \item \texttt{indyna3d} (\texttt{fd3d\_theo.f}) --- integra las ecuaciones
        elastodinámicas 3D por diferencias finitas con la condición de borde de
        fractura (fricción slip-weakening). Produce la historia de
        \textbf{slip-rate} $\dot{s}(i,k,t)$ sobre la falla.
  \item \texttt{seismograms} (\texttt{convm\_dynsource.f}) --- convoluciona el
        slip-rate con las funciones de Green de AXITRA y entrega los
        desplazamientos en las estaciones.
  \item \texttt{calcmisfit} --- compara con los datos observados (norma $L_2$
        rotada R/T/Z, ventanas P/S).
\end{enumerate}

El bucle de inversión (Neighbourhood Algorithm, \texttt{NA\_src/na.f}) llama
repetidamente a esta cadena buscando los parámetros que minimizan el misfit.

% ===========================================================================
\section{Ecuaciones elastodinámicas: formulación velocidad--esfuerzo}
% ===========================================================================

El medio es elástico, lineal e isótropo. Las incógnitas son el campo de
\textbf{velocidad de partícula} $\mathbf{v}=(v_x,v_y,v_z)$ y el \textbf{tensor de
esfuerzos} $\boldsymbol\sigma$. Las ecuaciones de gobierno son la conservación
del momento (segunda ley de Newton para un continuo)
\begin{equation}
  \rho\,\pd{v_i}{t} \;=\; \sum_j \pd{\sigma_{ij}}{x_j},
  \label{eq:momentum}
\end{equation}
y la ley constitutiva de Hooke en su forma de \emph{tasa} (derivada temporal),
\begin{equation}
  \pd{\sigma_{ij}}{t}
  \;=\; \lambda\,\delta_{ij}\,(\nabla\!\cdot\!\mathbf{v})
      \;+\; \mu\!\left(\pd{v_i}{x_j}+\pd{v_j}{x_i}\right),
  \label{eq:hooke}
\end{equation}
donde $\rho$ es la densidad y $\lambda,\mu$ las constantes de Lamé.

\paragraph{Correspondencia con el código.}
En el \texttt{common/disp/} de \texttt{parstat} los arreglos \texttt{u1,v1,w1}
\emph{son las velocidades} $v_x,v_y,v_z$ (aunque los comentarios hablen de
``displacement fields''): la actualización en \texttt{dvel} es exactamente
\eqref{eq:momentum}. Los seis arreglos \texttt{xx,yy,zz,xy,yz,xz} son las
componentes de $\boldsymbol\sigma$. Las propiedades del medio se guardan como
\begin{equation}
  \texttt{d1}=\rho,\qquad
  \texttt{mu1}=\mu=\rho\,v_s^2,\qquad
  \texttt{lam1}=\lambda=\rho\,(v_p^2-2v_s^2),
\end{equation}
que es lo que arma \texttt{rdmd3d} al leer \texttt{model.dat} (un modelo 1D por
capas, $v_p,v_s,\rho$, replicado en todo el plano horizontal).

% ===========================================================================
\section{Malla escalonada y operadores de diferencias finitas}
% ===========================================================================

\subsection{Malla escalonada (staggered grid)}
El esquema es el clásico de Madariaga (1976), popularizado por Virieux (1986):
las distintas componentes de $\mathbf{v}$ y $\boldsymbol\sigma$ viven en
posiciones desplazadas medio paso $\half\,\dd h$ entre sí, y el tiempo se avanza
en \emph{leapfrog} (las velocidades en $t+\half\dd t$, los esfuerzos en
$t+\dd t$). Esto da precisión de segundo orden en el tiempo con un solo nivel de
almacenamiento. El paso espacial es uniforme, \texttt{dh}, e isótropo.

\subsection{Operador espacial de 4.º orden}
Las derivadas espaciales usan el operador escalonado de cuarto orden, con
coeficientes (\texttt{deriv.f}: \texttt{c1=9./8.}, \texttt{c2=-1./24.})
\begin{equation}
  \pd{f}{x}\bigg|_{i}
  \;\approx\;
  \frac{1}{\dd h}\Big[\,c_1\big(f_{i+\frac12}-f_{i-\frac12}\big)
                     + c_2\big(f_{i+\frac32}-f_{i-\frac32}\big)\Big],
  \qquad c_1=\tfrac{9}{8},\; c_2=-\tfrac{1}{24}.
  \label{eq:fd4}
\end{equation}
En el código el medio-índice se materializa como diferencias entre nodos
vecinos; por ejemplo, la actualización de $v_x$ (\texttt{uxx1}) es
\begin{equation}
  v_x^{\,t+\frac12} = v_x^{\,t-\frac12}
   + \frac{\dd t}{\rho\,\dd h}\!\left[
       c_1(\sigma_{xx}^{i}-\sigma_{xx}^{i-1})+c_2(\sigma_{xx}^{i+1}-\sigma_{xx}^{i-2})
       + (\text{términos } xy,\,xz)\right],
\end{equation}
réplica directa de \eqref{eq:momentum} con el operador \eqref{eq:fd4}.
Análogamente \texttt{vyy1}, \texttt{wzz1} avanzan $v_y,v_z$, y \texttt{dstres}
avanza los esfuerzos según \eqref{eq:hooke} con $a=\lambda+2\mu$, $b=\lambda$
para las normales, y la media armónica/aritmética de $\mu$ entre nodos vecinos
(\texttt{xmu=xm1+xm2}) para las cizallas $\sigma_{xy},\sigma_{xz},\sigma_{yz}$.

\subsection{Orden reducido en los bordes}
En la primera capa de nodos junto a cada cara se usa un operador de
\textbf{segundo orden} (\texttt{uxx0/vyy0/wzz0} para velocidades vía
\texttt{bnd2d}; \texttt{sxx}, \texttt{sxy0/syz0/sxz0} para esfuerzos vía
\texttt{strbnd}), porque el operador de 4.º orden necesita dos nodos a cada lado.

\subsection{Estabilidad (CFL)}
La condición de estabilidad del esquema explícito es de tipo Courant. El código
la reporta como (\texttt{wrpm3d}, formato 70):
\begin{equation}
  \mathrm{CFL} \;=\; \frac{v_p^{\max}\,\dd t}{\dd h} \;<\; 0.5 .
  \label{eq:cfl}
\end{equation}
$v_p^{\max}$ es la mayor velocidad $P$ del modelo. Hay que elegir $\dd t$ acorde
a $\dd h$ y al modelo para no violar \eqref{eq:cfl}.

% ===========================================================================
\section{Condiciones de borde del dominio}
% ===========================================================================

\subsection{Borde absorbente (ABC)}
Para que la malla finita no refleje energía hacia adentro, \texttt{abc} aplica
una condición \textbf{paraxial de primer orden} (Clayton--Engquist) en las seis
caras. Para una onda que sale normal a una cara, la actualización tiene la forma
\begin{equation}
  f^{\text{borde}} \mathrel{+}= \frac{c\,\dd t}{\dd h}\,(f^{\text{interior}}-f^{\text{borde}}),
  \qquad
  c=\sqrt{\tfrac{\lambda+2\mu}{\rho}}=v_p \;\text{(normal)},\quad
  c=\sqrt{\tfrac{\mu}{\rho}}=v_s \;\text{(tangencial)},
\end{equation}
donde se usa $v_p$ para la componente normal a la cara y $v_s$ para las
tangenciales. Es una ABC sencilla (la actualización FD3D\_TSN la reemplaza por
PML, mucho más eficaz).

\subsection{Superficie libre}
Existen rutinas de superficie libre (\texttt{fres} para esfuerzos: $\sigma_{zz},
\sigma_{xz},\sigma_{yz}=0$ en $z=n_z$ vía imágenes antisimétricas; \texttt{fuvw}
para velocidades). En la configuración que nos interesa (fuente profunda,
$\sim$123 km) el bucle principal de \texttt{indyna3d} usa ABC en todas las caras;
la superficie libre del semiespacio real la aporta AXITRA al sintetizar.

% ===========================================================================
\section{La falla: geometría, esfuerzos y fricción}
% ===========================================================================

\subsection{Geometría de la falla en la malla}
La falla es un \textbf{plano} ubicado en $j=n_{ysc}$ (índice fijo en $y$); su
\textbf{normal es $\hat y$}. El deslizamiento (slip) es en dirección $z$, es
decir \textbf{dip-slip} (comentario en \texttt{fd3d\_theo.f}: ``modified
september 2011 for dip slip''). La tracción de cizalla relevante sobre el plano
es $\sigma_{yz}$, y la fricción se impone sobre \emph{esa} componente.

\paragraph{Modelo de media falla (antisimetría).} Sólo se simula la mitad del
medio: la velocidad de deslizamiento se toma como
$\dot{s}=2\,v_z(i,n_{ysc}{+}1,k)$ y el slip acumulado y el momento llevan el
factor $2$ (líneas \texttt{2.*w1}, \texttt{2.*slip} en \texttt{indyna3d} y
\texttt{wrss3d}). Esto explota la antisimetría del campo a través del plano de
falla.

\subsection{Esfuerzo inicial: \texttt{mkstress}}
\texttt{mkstress} llena el campo de tracción inicial de cizalla
$\sigma^0(i,j)\equiv\texttt{strinix}$ sobre el plano. Para cada punto $(i,j)$ se
evalúa su posición en el marco propio de la elipse (centro $(x_o,y_o)$, rotada un
ángulo $\phi$):
\begin{align}
  x' &= (i-x_o)\cos\phi + (j-y_o)\sin\phi, \\
  y' &= -(i-x_o)\sin\phi + (j-y_o)\cos\phi, \qquad
  r_e = \sqrt{\left(\tfrac{x'}{a}\right)^2+\left(\tfrac{y'}{b}\right)^2}.
\end{align}
El punto está \emph{dentro de la aspereza} si $r_e<1$. Entonces
\begin{equation}
  \sigma^0(i,j)=
  \begin{cases}
    \sigma_{\text{in}} \;(=\texttt{strbin}=T_e\cdot 10^6\ \text{Pa}), & r_e<1 \quad(\text{aspereza}),\\[2pt]
    \sigma_{\text{out}}\;(=-10^{8}\ \text{Pa}), & r_e\ge 1\quad(\text{fuera}).
  \end{cases}
\end{equation}
El valor exterior fuertemente negativo confina la ruptura a la aspereza (fuera no
puede romper). Además, dentro de un \textbf{círculo de nucleación} de radio $r$
centrado en el hipocentro $(x_a,y_a)=(n_x/2,n_y/2)$ se \emph{suma} un sobre-esfuerzo
$\sigma_{\text{nuc}}$ (\texttt{strain}) para forzar el inicio espontáneo:
\begin{equation}
  \text{si } (i-x_a)^2+(j-y_a)^2 \le r^2:\quad \sigma^0(i,j)\mathrel{+}= \sigma_{\text{nuc}}.
\end{equation}

\subsection{Esfuerzo de cedencia: \texttt{mkpeak}}
\texttt{mkpeak} arma el campo de \textbf{esfuerzo pico} (yield) $\tau_p(i,j)
\equiv\texttt{peak\_xz}$. En el código ambas ramas (dentro/fuera de la elipse)
asignan el mismo valor \texttt{strin}, de modo que el esfuerzo pico resulta
\textbf{uniforme} sobre la falla:
\begin{equation}
  \tau_p(i,j) = \tau_p \equiv \texttt{strin} = c_1\cdot \sigma_{\text{in}}
              = c_1\,T_e\cdot 10^6\ \text{Pa},
\end{equation}
donde $c_1=\texttt{rmodel(7)}$ (parámetro \texttt{cte1}). Es decir, $\tau_p$ es un
múltiplo del esfuerzo inicial de la aspereza.

\subsection{Ley de fricción slip-weakening y criterio de ruptura}
En cada paso de tiempo, tras actualizar los esfuerzos, se aplica la
\textbf{condición de borde de fractura} en el plano $j=n_{ysc}$
(\texttt{indyna3d}). Se trabaja con la tracción \emph{absoluta}
\begin{equation}
  \tau_{\text{abs}}(i,k) = \underbrace{\sigma_{yz}(i,n_{ysc},k)}_{\text{relativa, FD}}
                          + \underbrace{\sigma^0(i,k)}_{\text{inicial}} .
\end{equation}

\paragraph{(a) Criterio de iniciación.} Un punto aún intacto rompe cuando su
tracción absoluta supera el esfuerzo pico:
\begin{equation}
  \tau_{\text{abs}}(i,k) > \tau_p(i,k)
  \;\Longrightarrow\; \text{punto roto},\quad
  t_{\text{rup}}(i,k)=t .
\end{equation}

\paragraph{(b) Ley slip-weakening.} Una vez roto y deslizando, la tracción cae
linealmente con el slip acumulado $s$ hasta una distancia crítica $D_c$
(\emph{slip-weakening} clásico de Ida/Andrews; aquí en su versión Madariaga):
\begin{equation}
  \tau(i,k) = \tau_p(i,k)\cdot
  \begin{cases}
    1 - \dfrac{s}{D_{\max}}, & s \le D_{\max},\\[8pt]
    0, & s > D_{\max},
  \end{cases}
  \label{eq:sw}
\end{equation}
y se impone como condición de borde escribiendo en el campo relativo
$\sigma_{yz}(i,n_{ysc},k) = \tau(i,k) - \sigma^0(i,k)$. En el código
$D_{\max}=\texttt{rmodel(10)}$ es \textbf{$D_c/2$} (consistente con el factor 2 de
la media falla). El slip se integra explícitamente:
\begin{equation}
  s(i,k)\mathrel{+}= v_z(i,n_{ysc}{+}1,k)\,\dd t .
\end{equation}

\paragraph{(c) Cierre/no-penetración.} Si la apertura normal se vuelve negativa
($v_z(i,n_{ysc}{+}1,k)<0$, la falla tiende a cerrarse), el punto se ``des-rompe''
y se anulan las velocidades $v_x,v_z$ en el plano, evitando interpenetración.

\subsection{Salidas de la dinámica}
Al terminar el bucle temporal, \texttt{indyna3d} escribe (si \texttt{ioutput=1})
los mapas finales sobre la falla: tiempo de ruptura $t_{\text{rup}}$
(\texttt{ruptime.res}), slip final $2s$ (\texttt{slip.res}) y caída de esfuerzo
$\sigma_{yz}$ (\texttt{stressdrop.res}). El momento sísmico (sin escalar por
$\mu$ y área aún) se acumula como
\begin{equation}
  \texttt{amom} = \sum_{i,k} 2\,s(i,k).
\end{equation}

% ===========================================================================
\section{De slip-rate a sismogramas: \texttt{seismograms} + AXITRA}
% ===========================================================================

La radiación se calcula con el \textbf{teorema de representación}: el campo en
una estación es la convolución del slip-rate de cada subfalla con la función de
Green del medio 1D. Las Green las provee AXITRA (número de onda discreto,
Bouchon 1981) y vienen pre-calculadas en \texttt{axi.res} (leídas por
\texttt{readaxires} en los arreglos \texttt{uuxf,uuyf,uuzf} indexados por
frecuencia, fuente, receptor y las 6 dislocaciones elementales).

\subsection{Decimación de la fuente}
La malla dinámica (fina, $n_x\times n_y$) se \textbf{decima} por un factor
\texttt{idec=4} promediando bloques $4\times4$, de modo que la malla de fuentes
para la radiación es $4\times$ más gruesa que la del solver FD. Esto reduce el
número de fuentes puntuales en la convolución sin perder la cinemática gruesa.

\subsection{Tensor de momento por subfalla: \texttt{cmoment}}
Cada subfalla es un \textbf{doble par} (dislocación de cizalla) con orientación
$(\text{strike},\text{dip},\text{rake})$. \texttt{cmoment} convierte el mecanismo
en los seis coeficientes $a_{1..6}$ de las dislocaciones elementales que AXITRA
sabe propagar. El momento escalar por subfalla es $M_0=\mu\cdot \text{área}$
(el slip ya entra después vía el espectro de slip-rate), con
$\mu=\rho\,v_s^2$ evaluado en la capa donde cae la fuente. Los coeficientes son
las combinaciones estándar de $\sin/\cos$ de strike, dip y rake (ecuaciones
\texttt{x1..x6} $\to$ \texttt{a(1..6)}).

\subsection{Convolución en frecuencia y frecuencia compleja}
El slip-rate de cada subfalla $\dot{s}(i,j,t)$ se transforma a frecuencia con una
FFT. Para suprimir el \emph{wraparound} temporal (artefacto de la transformada
discreta) se usa una \textbf{frecuencia compleja} con parte imaginaria
\begin{equation}
  \omega = 2\pi f - i\,a_w/\text{(escala)} ,\qquad
  a_w = -\pi\,\texttt{aw}/t_l ,
\end{equation}
lo que equivale a multiplicar la señal por $e^{-a_w t}$ antes de transformar y
deshacerlo al final (técnica de Bouchon/Kennett). El desplazamiento espectral en
el receptor $r$ es la suma sobre subfuentes $s$ y sobre las 6 dislocaciones:
\begin{equation}
  u_c(\omega, r) \;=\; \sum_{s} \Big[\sum_{m=1}^{6} G^{c}_{m}(\omega,s,r)\,a_m(s)\Big]\,
                       \dot{S}(\omega,s),
  \qquad c\in\{x,y,z\},
\end{equation}
donde $\dot{S}(\omega,s)$ es el espectro del slip-rate de la subfalla $s$ y
$G^c_m$ las Green de AXITRA. Una FFT inversa devuelve las trazas en tiempo y se
deshace el peso exponencial de la frecuencia compleja.

\subsection{Filtrado de banda}
Cada traza sintética se filtra con un pasabanda Butterworth (\texttt{XAPIIR},
tipo \texttt{BU}, \texttt{BP}) entre \texttt{fr1} y \texttt{fr2} (típicamente
$0.02$--$0.1$ Hz). El resultado se vuelca a \texttt{predicted\_data} en el orden
$(x,y,z)$ por estación, que es el que espera el misfit. Si \texttt{ioutput=1}
también se escriben \texttt{dyn\_disp\_x/y/z}.

% ===========================================================================
\section{Función de misfit: \texttt{calcmisfit}}
% ===========================================================================

El misfit es una \textbf{norma $L_2$ normalizada}, calculada sobre componentes
\textbf{rotadas} radial/transversal/vertical (R/T/Z) y \textbf{ventaneada} por
fases. Es idéntico en filosofía al misfit cinemático de KDEllipsPy.

\subsection{Rotación a R/T/Z}
Con el acimut $\phi$ de cada estación (de \texttt{azi\_times.txt}), las
componentes N ($x$) y E ($y$) se rotan:
\begin{equation}
  R = N\cos\phi + E\sin\phi,\qquad
  T = E\cos\phi - N\sin\phi,\qquad
  Z = Z .
\end{equation}

\subsection{Ventanas por fase}
\begin{itemize}
  \item Onda \textbf{P}: se evalúan las componentes \textbf{radial} y
        \textbf{vertical} en una ventana que arranca en el tiempo de arribo
        $t_P$ y dura \texttt{time\_window} $=20$ s.
  \item Onda \textbf{S(H)}: se evalúa la componente \textbf{transversal} en una
        ventana que arranca en $t_S$ y dura $20$ s.
\end{itemize}
El muestreo de trabajo del misfit es $4$ Hz, así que el índice de inicio es
$\texttt{nint}(t_{P/S})\times 4$.

\subsection{Definición del misfit}
Acumulando sobre componentes, estaciones y muestras dentro de las ventanas:
\begin{equation}
  \text{misfit}
  \;=\; \frac{\displaystyle\sum \big(d^{\text{obs}}-d^{\text{pred}}\big)^2}
             {\displaystyle\sum \big(d^{\text{obs}}\big)^2}
  \;=\; \frac{\texttt{num}}{\texttt{den}} .
  \label{eq:misfit}
\end{equation}
Si la ruptura \textbf{no nuclea}, las sintéticas son nulas, $\texttt{num}=
\texttt{den}$ y el misfit da exactamente $1.0$; el código lo detecta y lo manda a
$\infty$ ($\texttt{misfitval}/0$) para que el NA descarte ese modelo.

% ===========================================================================
\section{Parámetros invertidos y algoritmo de inversión}
% ===========================================================================

El forward depende de \textbf{10 parámetros dinámicos} (\texttt{rmodel(1..10)} en
\texttt{forward\_modelling.f}; rangos en \texttt{input\_dynamic\_ellipse.txt}):

\begin{table}[h]
\centering
\begin{tabular}{>{\ttfamily}c l l}
\toprule
\normalfont\# & Símbolo / significado & Unidad \\
\midrule
rmodel(1)  & $a$ --- semieje 1 de la aspereza elíptica      & puntos de malla \\
rmodel(2)  & $b$ --- semieje 2 de la aspereza               & puntos de malla \\
rmodel(3)  & $x_o$ --- centro de la elipse (strike)         & puntos de malla \\
rmodel(4)  & $y_o$ --- centro de la elipse (dip)            & puntos de malla \\
rmodel(5)  & $\phi$ --- ángulo de rotación de la elipse     & radianes \\
rmodel(6)  & $T_e$ --- esfuerzo inicial de la aspereza ($\sigma_{\text{in}}=T_e\,10^6$) & MPa \\
rmodel(7)  & $c_1$ --- $\tau_p = c_1\,\sigma_{\text{in}}$ (esfuerzo pico)   & adim. \\
rmodel(8)  & $c_2$ --- $\sigma_{\text{nuc}} = c_2\,\tau_p$ (sobre-esfuerzo de nucleación) & adim. \\
rmodel(9)  & $r$ --- radio del círculo de nucleación        & puntos de malla \\
rmodel(10) & $D_{\max}=D_c/2$ --- distancia crítica slip-weakening & m \\
\bottomrule
\end{tabular}
\caption{Los 10 parámetros del modelo dinámico. El hipocentro está fijo en el
centro de la malla $(x_a,y_a)=(n_x/2,n_y/2)$; \texttt{strbout}$=-10^8$ Pa y
\texttt{straout}$=0$ están \emph{hardcodeados}.}
\end{table}

\paragraph{Neighbourhood Algorithm (NA).} El bucle externo
(\texttt{Source/NA\_src/na.f}; Sambridge 1999) muestrea el espacio de 10
parámetros, evalúa el misfit \eqref{eq:misfit} de cada modelo vía la cadena
forward, y re-muestrea las celdas de Voronoi de menor misfit. Los controles
(número de iteraciones, tamaño de muestra inicial y por iteración, celdas a
re-muestrear, número de núcleos) salen de
\texttt{input\_dynamic\_ellipse.txt}. \texttt{fd3d\_direct} corre un único
forward del mejor modelo con \texttt{ioutput=1} para volcar todas las salidas.

% ===========================================================================
\section{Resumen de constantes y archivos clave}
% ===========================================================================

\begin{itemize}
  \item \textbf{Malla FD} (\texttt{parstat}): \texttt{nx=ny=nz=160} (en
        compilación). La malla de fuentes radiantes es $n_x/4 \times n_y/4$.
  \item \textbf{Operador 4.º orden}: $c_1=9/8$, $c_2=-1/24$.
  \item \textbf{Estabilidad}: $v_p^{\max}\dd t/\dd h < 0.5$.
  \item \textbf{Fricción}: slip-weakening sobre $\sigma_{yz}$, $\tau_p$ uniforme,
        $D_{\max}=D_c/2$.
  \item \textbf{Green}: AXITRA 1D (número de onda discreto), 6 dislocaciones
        elementales, frecuencia compleja $a_w$ para anti-wraparound.
  \item \textbf{Banda}: pasabanda Butterworth \texttt{XAPIIR} $[fr_1,fr_2]$.
  \item \textbf{Misfit}: $L_2$ normalizado, R/T/Z, ventanas P (R,Z) y S (T) de
        20 s desde $t_P/t_S$, muestreo 4 Hz.
  \item \textbf{I/O de la falla}: \texttt{stressin.dat}, \texttt{peakin.dat},
        \texttt{ruptime.res}, \texttt{slip.res}, \texttt{stressdrop.res},
        \texttt{sliprate.res}, \texttt{synth.res}, \texttt{dyn\_disp\_\{x,y,z\}}.
\end{itemize}

% ===========================================================================
\section{Implicancias para la reimplementación en Python}
% ===========================================================================

Para portar el \emph{forward} a Python conviene separar tres bloques:
\begin{enumerate}
  \item \textbf{Generador elipse $\to$ campos $(\sigma^0,\tau_p)$}: equivale a
        \texttt{mkstress}+\texttt{mkpeak}; es álgebra vectorizada trivial en
        NumPy y reutiliza la misma parametrización elíptica del lado cinemático.
  \item \textbf{Solver FD3D} (\texttt{indyna3d}): es el núcleo costoso. Opciones:
        (a) envolver el Fortran existente con \texttt{f2py} (igual que ya se hizo
        con AXITRA), o (b) adoptar \textbf{FD3D\_TSN} (Premus--Gallovič), que es
        el descendiente directo del mismo solver de Madariaga con
        \emph{traction-at-split-node}, PML y aceleración GPU. La salida que se
        necesita es la historia de slip-rate $\dot{s}(i,k,t)$.
  \item \textbf{Radiación + misfit}: la convolución con AXITRA ya está envuelta
        en \texttt{kdellipspy/axitra}, y el misfit \eqref{eq:misfit} coincide con
        el cinemático ya implementado. El NA en Python (\texttt{NAInversionModel})
        puede orquestar todo, sustituyendo el frágil \texttt{bash+sed} del legacy.
\end{enumerate}

\end{document}
